\section{Introduction}
\label{sec:introduction}

The Yang-Mills existence and mass gap problem, one of the seven Millennium Prize Problems, asks for a rigorous mathematical proof that non-Abelian gauge theories in four dimensions possess a mass gap—a strictly positive lower bound on the energy spectrum above the vacuum state.

\subsection{Scope and status of results}

This document is written as a \emph{rigorous foundations-and-roadmap} manuscript.
When the Yang--Mills Millennium problem is discussed, we distinguish carefully between:
\begin{itemize}
\item \textbf{Proved statements} (with complete proofs in the text),
\item \textbf{Conditional theorems} (clearly listing hypotheses that are not proved here), and
\item \textbf{Programmatic steps / conjectures} (motivated directions that remain open).
\end{itemize}

In particular, we do \emph{not} assert an unconditional construction of continuum four-dimensional Yang--Mills theory with a nonzero mass gap in this version. Instead we aim to (i) make the verified parts of the argument fully rigorous, and (ii) isolate a small set of precise, checkable open inputs on which a complete solution would depend.

\subsection{Architecture of the program}

The manuscript is organized into four stages that separate what is standard/verified from what requires genuinely new input.

\subsubsection*{Stage I: The Lattice Foundation (Finite Volume)}
We begin with the standard Wilson lattice regularization. Using the transfer matrix formalism and reflection positivity, we prove the existence of a spectral gap $\Delta_L(\beta) > 0$ for any finite lattice volume $L^3$ and any coupling $\beta$. The critical challenge is to prove that $\lim_{L \to \infty} \Delta_L > 0$ and that this gap survives the continuum limit $a \to 0$.

\subsubsection*{Stage II: Strong Coupling Control ($\beta < \beta_c$)}
In the strong coupling regime, we utilize the \textbf{Cluster Expansion} technique. We prove that the polymer expansion of the partition function converges absolutely (Theorem R.150.6). This establishes the exponential decay of correlations and a mass gap $\Delta \propto \ln(1/\beta)$, providing a rigorous anchor for the theory.

\subsubsection*{Stage III: The Uniformity Bridge (Intermediate Coupling)}
To bridge the gap between strong and weak coupling, one needs a uniform-in-volume mechanism that prevents the infinite-volume gap from collapsing at intermediate coupling. In this manuscript, proposed mechanisms are formulated precisely and labeled as conditional where appropriate:

\begin{itemize}
\item \textbf{Adjoint interpolation (conditional idea):} Adding massive adjoint fermions relates pure YM to supersymmetric regimes. Center symmetry alone prevents \emph{deconfinement} transitions, but does not exclude unrelated bulk transitions; any analyticity/no-phase-transition claim is therefore stated conditionally and requires a separate proof.
\item \textbf{Functional-inequality approach (conditional idea):} A uniform log-Sobolev/spectral-gap bound would provide a robust route to infinite-volume control. Establishing such uniformity for non-Abelian lattice gauge measures at all couplings is itself a major open analytic step; the manuscript isolates the exact inequality needed.
\end{itemize}

\subsubsection*{Stage IV: The Continuum Limit ($\beta \to \infty$)}
We discuss the continuum limit $\beta \to \infty$ and the scale-setting problem, emphasizing which ingredients are rigorous (e.g. reflection positivity on the lattice, known RG constructions in restricted settings) and which inputs remain open for four-dimensional non-Abelian pure YM.

\begin{itemize}
\item \textbf{String tension and confinement:} Positivity of a string tension for all couplings is not asserted as a theorem here unless proved with correct non-Abelian correlation inequalities.
\item \textbf{Spectral relations:} Phenomenological relations between $\Delta$ and $\sigma$ (sometimes called Giles--Teper-type relations) are treated as motivation and clearly labeled as non-rigorous unless proved.
\item \textbf{Continuum limit:} Mosco/Dirichlet-form convergence is described as a framework; the required uniform estimates are stated explicitly.
\end{itemize}

\subsection{How obstacles are handled}
Rather than a ``patch list'', we organize technical issues as \emph{definitions, lemmas, and explicit open inputs}.
When an ingredient is standard in constructive QFT (e.g. reflection positivity on the lattice), we provide a complete proof or a precise reference path.
When an ingredient is genuinely open in non-Abelian 4D YM (e.g. global analyticity/no-bulk-transition statements, uniform correlation inequalities, or unconditional nonperturbative bounds at all couplings), we state it as an open problem/hypothesis and track its logical dependencies.


